\chapter{Metodologia}\label{cap:metodologia}

Este capítulo detalha as modificações desenvolvidas nos módulos de identificação do
\emph{pipeline} do desafio SoccerNet-GSR, mais especificamente no reconhecimento dos
números das camisolas e na atribuição da equipa e do lado do campo. Começa-se por enquadrar
estas intervenções na arquitetura geral da \emph{baseline}, deixando claro o que foi alterado
e o que se manteve intacto. Em seguida, apresentam-se individualmente o novo módulo de leitura
visual e as correções aplicadas à classificação de equipas. Por fim, descrevem-se o ambiente
experimental e o protocolo de avaliação sobre os quais assentam os resultados do capítulo
seguinte.


\section{Visão Geral do Sistema}

A arquitetura da \emph{baseline} está ilustrada na Figura~\ref{fig:gsr-pipeline}. Dos módulos
que a compõem, este trabalho intervém unicamente na etapa final de identificação, aquela que
atribui a cada \emph{tracklet} a equipa, a função e o número de camisola.

As alterações concentram-se em dois dos seus componentes. No primeiro, a leitura do número de
camisola, o módulo de OCR da \emph{baseline} deu lugar a um modelo multimodal de visão e
linguagem. No segundo, a atribuição de equipa e do lado do campo, corrigiram-se erros de
agrupamento e de mapeamento e recalibraram-se os limiares de decisão. Todos os restantes
módulos, a deteção, a reidentificação, o \emph{tracking} e a calibração, permaneceram intocados.

A escolha destes dois componentes não é arbitrária. São os que menos atenção têm recebido na
literatura, como evidencia a revisão do estado da arte, e aqueles em que a solução de referência
revela maior margem de melhoria, de forma particularmente acentuada no reconhecimento do número de
camisola. Isto torna-os o ponto de intervenção com maior retorno esperado.

Concentrar o esforço apenas nestes dois módulos, em vez de o dispersar por todo o \emph{pipeline},
permitiu ainda estudá-los com uma profundidade que uma intervenção mais alargada dificilmente
comportaria, em linha com o âmbito de uma dissertação de mestrado. Isolá-los, mantendo tudo o resto
inalterado, garante que qualquer variação no desempenho final resulte diretamente das melhorias
introduzidas, e não de uma reotimização global. Por serem independentes entre si, as duas
alterações integram-se no \emph{pipeline} sem exigir o reentreino das etapas anteriores, mantendo a
compatibilidade direta com a \emph{baseline} pública disponibilizada por \citet{somers2024soccernet}.

\section{Reconhecimento de Número de Camisola}\label{sec:met-jnr}

Na solução de referência, o reconhecimento dos números das camisolas baseia-se na biblioteca
MMOCR, que aplica um \emph{pipeline} genérico de deteção e reconhecimento de texto sobre o recorte de
cada jogador. Por atuar estritamente ao nível do carácter e sem incorporar conhecimento
\emph{a priori} sobre o domínio, esta abordagem trata o número como texto arbitrário e ignora
duas propriedades fundamentais do cenário desportivo.

A primeira é de cariz espacial. O número encontra-se numa região anatomicamente delimitada, a
zona dorsal do tronco, cuja localização pode ser estimada com precisão por um modelo de pose, em
vez de ser procurada de forma exaustiva em toda a imagem. A segunda é de cariz temporal. O mesmo
jogador é observado ao longo de múltiplos \emph{frames} consecutivos, e esta redundância,
frequentemente desaproveitada pelos métodos de inferência por \emph{frame} isolado, revela-se
crítica em situações de baixa resolução, oclusão parcial ou poses desfavoráveis.

Para tirar partido destas especificidades, o módulo proposto reorganiza a leitura numa sequência
de três fases interdependentes (Figura~\ref{fig:jnr-flow}). Inicialmente, um modelo de estimativa
de pose (YOLO26m-pose) extrai os pontos de interesse (\emph{keypoints}) do tronco e isola a
região dorsal, eliminando o fundo e as áreas irrelevantes do recorte. Em seguida, essa região é
processada por um modelo multimodal de visão e linguagem (\emph{Vision-Language Model}, VLM), o
Qwen2.5-VL, capaz de inferir conjuntamente sobre a informação visual e o contexto semântico do
equipamento e de distinguir o número de outros elementos gráficos ou patrocínios. Por fim, as
predições individuais associadas a cada \emph{tracklet} são consolidadas por um mecanismo de
votação temporal, convertendo uma sequência de inferências instantâneas numa decisão final mais
robusta e consistente.

\begin{figure}[htbp]
    \centering
    \includegraphics[width=\textwidth]{jnr_flow.png}
    \caption{Fluxo do módulo de reconhecimento do número de camisola. A estimativa de pose
    delimita a região dorsal, cujo recorte é interpretado pelo modelo multimodal Qwen2.5-VL
    quantizado a 4 bits. As leituras sucessivas de cada \emph{tracklet} são consolidadas por
    votação temporal numa decisão final única.}
    \label{fig:jnr-flow}
\end{figure}

\subsection{Recorte da Região Dorsal por Estimativa de Pose}

Para concentrar a leitura na zona onde o número efetivamente aparece, cada deteção é primeiro
processada por um modelo de estimativa de pose, o YOLO26m-pose. A partir dos \emph{keypoints} do
tronco, define-se uma região de interesse sobre as costas do jogador, alargada lateralmente para
acomodar variações de postura e enquadramento. Este recorte elimina o fundo, os membros e a zona
frontal, entregando ao modelo de leitura apenas o conteúdo relevante.

Antes de prosseguir, cada recorte passa por um filtro de qualidade que descarta os casos sem
informação utilizável. São rejeitados os recortes com dimensões inferiores a um limiar mínimo,
tipicamente jogadores muito distantes da câmara, e os recortes desfocados, quantificados pela
variância do Laplaciano. Este passo é determinante, uma vez que uma imagem sem resolução
suficiente não contém o número a ler e a sua submissão ao modelo apenas introduziria ruído. Ao
filtrar à entrada, garante-se que o modelo de leitura recebe somente recortes com hipótese real
de sucesso.

\subsection{Leitura por Modelo Multimodal de Visão e Linguagem}

O núcleo deste módulo assenta no modelo multimodal Qwen2.5-VL-3B-Instruct, responsável por
inferir e reportar o número a partir dos recortes da região dorsal de cada jogador. A opção por
uma arquitetura de visão e linguagem (\emph{Vision-Language Model}, VLM), em detrimento de um
motor de OCR dedicado, fundamenta-se na sua capacidade de interpretação contextual, que lhe
permite distinguir a numeração oficial de outros elementos gráficos presentes no equipamento,
como patrocínios, logótipos de marcas ou textos avulsos, que de outro modo introduziriam ruído.

A principal limitação prática derivou da infraestrutura computacional disponível, uma GPU móvel
equipada com apenas 6~GB de VRAM. Para viabilizar a execução dentro deste limite de memória, o
modelo foi carregado com quantização a 4 bits (\emph{4-bit quantization}), o que reduz a sua
ocupação de VRAM para aproximadamente 2~GB e assegura a sua coexistência com os restantes
módulos do \emph{pipeline}. Adicionalmente, para prevenir exceções de memória (\emph{out-of-memory})
durante as fases subsequentes de calibração, a instância do modelo é explicitamente libertada da
VRAM no término do processamento de cada vídeo.

A inferência do modelo foi moldada por engenharia de \emph{prompt}, assente em duas diretrizes
fundamentais:

\begin{enumerate}
    \item \textbf{Integridade estrutural da numeração.} O modelo é instruído a extrair a
    totalidade dos dígitos, prevenindo a omissão de um dígito inicial parcialmente ocluído, uma
    falha recorrente que tende a truncar números compostos, por exemplo classificando um 14 como
    4.
    \item \textbf{Política conservadora de recusa.} O modelo deve emitir uma resposta nula
    sempre que a amostra apresente baixa resolução, desfocagem acentuada ou ambiguidade visual,
    em vez de arriscar uma leitura incerta.
\end{enumerate}

Esta estratégia conservadora reflete uma decisão de projeto importante. No contexto da
reconstrução do estado de jogo e da métrica oficial do SoccerNet-GSR, a propagação de um erro de
classificação é significativamente mais penalizadora do que uma ausência temporária de leitura.
Por fim, para otimizar a eficiência das chamadas ao modelo, vários recortes do mesmo
\emph{tracklet} são apresentados na mesma \emph{prompt}. Este agrupamento é, contudo, limitado a
um número reduzido de \emph{tracklets} em cada chamada, de forma a respeitar o limiar a partir do
qual o mecanismo de atenção do modelo se satura e a precisão das respostas se degrada.

\subsection{Agregação Temporal por Votação}

A leitura efetuada sobre um recorte isolado é, por natureza, ruidosa. Uma pose desfavorável, uma
oclusão momentânea ou variações de iluminação são suficientes para induzir erros na inferência de
um \emph{frame}, mesmo quando a numeração se mantém globalmente legível ao longo do
\emph{tracklet} (Figura~\ref{fig:jnr-crops}). Para converter este conjunto de leituras individuais
numa decisão única e fiável, adotou-se um esquema de votação temporal estruturado em duas fases
hierárquicas, concebido para mitigar o ruído em diferentes escalas de tempo.

\begin{figure}[htbp]
    \centering
    \includegraphics[width=\textwidth]{jnr_crops_tracklet71.png}
    \caption{Recortes da região dorsal de um mesmo \emph{tracklet} ao longo do tempo. O número do
    jogador é nítido apenas em parte dos \emph{frames}, sendo comprometido nos restantes por
    desfoque, oclusão ou pose desfavorável. É esta variabilidade que justifica tanto o filtro de
    qualidade à entrada como a consolidação das leituras por votação temporal.}
    \label{fig:jnr-crops}
\end{figure}

Na primeira fase, os recortes de cada \emph{tracklet} são particionados em janelas temporais.
Cada janela apenas emite uma predição válida se uma fração dominante das suas leituras convergir
no mesmo valor, filtrando assim variações momentâneas e descartando instâncias erráticas entre
\emph{frames} adjacentes. Na segunda fase, a numeração final atribuída ao \emph{tracklet}
corresponde à hipótese mais votada entre as janelas, desde que cumpra uma quota mínima do total
emitido. Deste modo, a primeira fase filtra o ruído de alta frequência ao nível do \emph{frame},
enquanto a segunda arbitra leituras concorrentes que persistem no tempo, conferindo robustez ao
processo.

Um aspeto determinante deste mecanismo é a inclusão das leituras nulas no apuramento final, sendo
estas contabilizadas negativamente contra qualquer número candidato. Ao integrar as janelas sem
leitura no denominador da votação, impede-se que uma leitura incorreta e isolada, suportada por
uma única janela válida entre várias vazias, obtenha suporte suficiente para prevalecer. Esta
abordagem neutraliza o risco de falsos positivos decorrente do critério conservador do modelo de
visão, que gera deliberadamente uma elevada proporção de respostas nulas. A exigência de confiança
na inferência individual e a penalização da incerteza na agregação operam, assim, de forma
concertada, assegurando que a decisão final priorize a consistência temporal em detrimento de
evidências pontuais.

\subsection{Alternativas Consideradas}

A solução final para a leitura do número não foi imediata. Resultou da comparação de várias famílias
de OCR e de modelos de visão e linguagem, assim como de estratégias de pré-processamento da imagem.
A Tabela~\ref{tab:alternativas} resume as abordagens exploradas e a razão do seu descarte.

\begin{table}[htbp]
\centering
\caption{Alternativas de leitura do número consideradas e razão do seu descarte.}
\label{tab:alternativas}
\renewcommand{\arraystretch}{1.35}
\setlength{\tabcolsep}{6pt}
\begin{tabular}{|p{3.6cm}|p{9.0cm}|}
\hline
\textbf{Abordagem} & \textbf{Razão do descarte} \\
\hline
EasyOCR e OCR clássico &
Mesma limitação estrutural do MMOCR, falha em números pequenos ou inclinados e não
distingue o número de outros elementos gráficos. \\
\hline
PaddleOCR (com \emph{fine-tuning}) &
Falha de segmentação (\emph{segfault}) ao inicializar no mesmo processo que o CUDA do
\emph{pipeline}, e mesmo isolado não recuperou leituras face ao Qwen. \\
\hline
HunyuanOCR &
\emph{VLM} mais pesado, sem ganho de exatidão que compensasse o custo de VRAM e de
latência no \emph{hardware} disponível. \\
\hline
Qwen2.5-VL-7B &
Requer cerca de 10\,GB de VRAM, acima dos 6\,GB da GPU, o que inviabiliza a sua
utilização. \\
\hline
Super-resolução (RealESRGAN) &
A informação em falta não existe na imagem de origem, e a ampliação amplifica o ruído
em vez de a recuperar. \\
\hline
SAM-3D-Body &
Reconstrução tridimensional do corpo demasiado pesada para correr por \emph{frame} e
desadequada a um simples sinal de orientação. \\
\hline
Classificador de orientação &
O ganho de filtrar as \emph{frames} sem as costas visíveis não compensou, pois a votação
por \emph{tracklet} já absorve naturalmente essas \emph{frames}. \\
\hline
\end{tabular}
\end{table}


Estas experiências convergem numa conclusão. Em condições de baixa resolução, o fator decisivo não é
a escolha do motor de OCR mais potente, mas sim entregar ao modelo a imagem certa, o recorte da
região dorsal obtido por pose, e agregar bem as leituras ao longo do tempo. O Qwen2.5-VL-3B foi
escolhido pelo equilíbrio entre exatidão e ocupação de memória, e não por ser o modelo de maior
dimensão.

\section{Atribuição de Equipa e do Lado do Campo}

Na solução de referência, a atribuição de equipa resulta de um agrupamento (\emph{clustering}) por
\emph{k-means} aplicado às características de reidentificação que o PRTreID extrai de cada jogador. O
\emph{k-means} é um algoritmo de agrupamento não supervisionado que reparte as observações num
número fixo de grupos, aqui dois, um por equipa, atribuindo cada observação ao grupo cujo centro lhe
está mais próximo e ajustando esses centros de forma iterativa até a partição estabilizar. Aplicado
a estes vetores de características (\emph{features}), o método reúne num mesmo grupo os jogadores que
ficam próximos no espaço latente, o que tende a corresponder aos que envergam o mesmo equipamento.
Os \emph{tracklets} são assim separados em dois grupos e cada grupo é depois associado a um dos lados
do campo, a partir da posição dos jogadores no relvado. As alterações introduzidas incidem sobre os
dois momentos deste processo, o agrupamento propriamente dito e o mapeamento dos grupos para os
lados, e incluem ainda a recalibração dos limiares de decisão que os governam.


\subsection{Agrupamento por Características Visuais}

A solução de referência atribui as equipas num único passo. Aplica o \emph{k-means} com dois grupos
sobre a média, por \emph{tracklet}, das características de reidentificação extraídas pelo PRTreID, e
assume que os dois grupos obtidos correspondem às duas equipas. Isto funciona bem quando os
equipamentos têm cores claramente diferentes, mas falha quando as duas equipas usam cores parecidas,
porque nesse caso as características de reidentificação deixam de as distinguir.

A alteração introduzida consiste em aferir a qualidade deste agrupamento antes de o aceitar e em
dispor de uma alternativa para quando ele não é de confiança. O agrupamento por reidentificação é
sempre calculado primeiro, por ser o sinal mais informativo quando funciona, e é depois submetido a
dois testes independentes. Basta que um deles acuse falta de fiabilidade para que a alternativa
entre em ação.

O primeiro teste recorre ao coeficiente de silhueta, uma medida que avalia se os dois grupos ficaram
bem definidos. Para cada \emph{tracklet}, compara a distância aos elementos do seu próprio grupo com
a distância aos do outro grupo, e resume essa relação num valor entre $-1$ e $1$. Valores próximos de
$1$ indicam grupos coesos e bem separados, ao passo que valores baixos indicam grupos que se
confundem. Quando a média deste valor sobre todos os \emph{tracklets} fica abaixo de $0.28$,
considera-se que a reidentificação não separou as duas equipas.

O segundo teste dirige-se a uma situação que a silhueta não capta. Quando a jogada decorre
concentrada numa das metades do campo, com pelo menos $85\%$ dos \emph{tracklets} do mesmo lado, e
os centros dos dois grupos estão simultaneamente próximos no espaço de reidentificação, a uma
distância euclidiana inferior a $2.04$, a separação obtida é tomada como frágil ainda que a silhueta
não a denuncie. As posições manifestamente fora do relvado, resultantes de falhas pontuais de
calibração, são descartadas antes desta contagem.

Verificada qualquer uma das duas condições, passa-se a agrupar pela cor da camisola em vez das
características de reidentificação. Essa cor não é medida neste momento, mas sim um descritor que o
sistema já calcula para cada jogador na fase de reidentificação, a partir da zona do tronco isolada
pela máscara de segmentação do corpo, e que resume a cor do equipamento sob a forma de um
histograma. Aplica-se novamente o \emph{k-means}, agora apenas sobre este descritor de cor. Usa-se
isolado, e não combinado com o vetor de reidentificação, porque os seus poucos valores seriam
abafados pelas centenas de dimensões da reidentificação e a cor deixaria de pesar na separação. A
Figura~\ref{fig:team_cluster} mostra, para um vídeo real, a separação dos dois grupos tanto no espaço
de reidentificação como no descritor de cor.

A aceitação desta alternativa não é incondicional, mas a exigência difere consoante o teste que a
desencadeou. Quando foi o coeficiente de silhueta, requer-se que a divisão por cor seja equilibrada,
com o grupo menor a reunir pelo menos $20\%$ dos \emph{tracklets}, mantendo-se o agrupamento
original em caso contrário, para não trocar uma separação fraca por outra ainda menos fiável.
Quando foi a segunda condição, a alternativa é aceite sem essa exigência, por se ter optado por
preservar o comportamento original desta regra, cuja alteração introduzia regressões no conjunto de
teste.

\begin{figure}[htbp]
    \centering
    \includegraphics[width=\textwidth]{team_cluster_space.png}
    \caption{Agrupamento dos \emph{tracklets} de um vídeo nas duas equipas. À esquerda, a projeção
    das características de reidentificação em duas dimensões, onde o \emph{k-means} separa os dois
    grupos com um coeficiente de silhueta de $0.42$. À direita, o descritor de cor da camisola dos
    mesmos \emph{tracklets}, que preserva a separação, com silhueta de $0.32$, e serve de alternativa
    quando a reidentificação é pouco fiável.}
    \label{fig:team_cluster}
\end{figure}

\subsection{Mapeamento para os Lados do Campo}\label{sec:mapeamento}

Definidos os dois grupos, resta associar cada um ao seu lado do campo. Esta decisão apoia-se na
posição dos jogadores no relvado, projetada pela calibração, e resolve-se por uma cascata de quatro
critérios de exigência decrescente.

O primeiro é uma votação por \emph{tracklet}. Cada \emph{tracklet} vota no lado correspondente à
sua posição média, depois de filtrados os valores atípicos, e um grupo é atribuído a um lado
quando uma fração dominante dos seus \emph{tracklets} aponta nesse sentido. Quando esta votação é
inconclusiva, por os dois grupos surgirem ambíguos ou do mesmo lado, entra um segundo critério que
compara a tendência relativa dos dois grupos, atribuindo o lado direito àquele que reúne
proporcionalmente mais \emph{tracklets} desse lado, ainda que nenhum atinja a maioria absoluta. Se
ambas as votações falharem, comparam-se diretamente as posições médias dos dois grupos e o lado
direito é dado ao grupo cuja média está mais à direita. Só quando essas médias são demasiado
próximas para constituírem sinal fiável, tipicamente quando toda a jogada decorre concentrada numa
zona do campo, se recorre a um último desempate. Nesse caso, toma-se como equipa atacante aquela
que, ao longo da sequência, mais avança em direção à baliza defendida pelo guarda-redes, ficando a
outra associada ao lado oposto. A Figura~\ref{fig:team_side} resume esta cascata de decisão.

Os limiares que governam esta cascata, tanto a maioria exigida na votação como as margens mínimas
dos restantes critérios, foram recalibrados em conjunto sobre os dados de teste e de validação. A
recalibração obedeceu a uma condição estrita, a de não degradar nenhum vídeo que já estivesse
correto, de modo que cada ajuste só foi aceite quando resolvia casos falhados sem introduzir novas
falhas.

\begin{figure}[htbp]
    \centering
    \includegraphics[width=0.95\textwidth]{team_side_flow.png}
    \caption{Fluxograma do processo que atribui cada grupo ao seu lado do campo, aplicando os
    critérios por ordem até um deles resolver a atribuição.}
    \label{fig:team_side}
\end{figure}

\FloatBarrier

\section{Ambiente Experimental}

Esta secção descreve as condições em que o sistema foi desenvolvido e avaliado. Apresenta o ambiente
computacional e o conjunto de dados, enumera os hiperparâmetros que governam os módulos
intervencionados e define o protocolo de avaliação sobre o qual assentam os resultados do capítulo
seguinte.

\subsection{Ambiente Computacional e Conjunto de Dados}

Todo o trabalho foi implementado com base no \emph{framework} TrackLab, responsável pela gestão e
execução sequencial dos módulos do \emph{pipeline}, tendo como suporte a biblioteca PyTorch. O
desenvolvimento decorreu em Python~3.10 e a execução final sobre o conjunto de treino em
Python~3.12, com as versões das bibliotecas fixadas e sem qualquer alteração ao código dos módulos.
A principal restrição computacional decorreu da infraestrutura de \emph{hardware}
disponível, um sistema portátil equipado com uma GPU NVIDIA RTX~3060 com apenas 6~GB de VRAM. Esta
limitação condicionou várias das decisões metodológicas adotadas, com destaque para a quantização a
4 bits do modelo de visão e linguagem, com recurso à biblioteca \emph{bitsandbytes}, e a desalocação
manual da memória de vídeo após o processamento de cada sequência.

A validação experimental utilizou o conjunto de dados SoccerNet-GSR, constituído por
sequências de vídeo anotadas e divididas nos subconjuntos de treino, validação e teste.
O protocolo de avaliação abrangeu a totalidade das 160 sequências anotadas,
correspondentes aos conjuntos de treino (57 vídeos), de validação (54 vídeos) e de teste
(49 vídeos). Para cada interveniente, este conjunto disponibiliza a anotação de
referência (\emph{ground truth}) da posição no terreno de jogo, da equipa, da função e do
número de camisola, permitindo tanto a avaliação integrada do sistema como o teste
isolado de cada módulo.

O papel de cada subconjunto neste trabalho difere do habitual e convém precisá-lo. Nenhum dos
módulos desenvolvidos envolve treino de modelos, pelo que o subconjunto de treino não foi usado
para aprender parâmetros. Os limiares de decisão foram afinados sobre os conjuntos de teste e de
validação, como se descreve na Secção~\ref{sec:mapeamento}, e o conjunto de treino permaneceu
inteiramente por observar até à fase final. Serve, por isso, o propósito de um conjunto de
retenção, ou seja, de medir o desempenho em dados que não intervieram em nenhuma decisão de
conceção.

\subsection{Parametrização dos Módulos}

O comportamento dos módulos desenvolvidos é definido por um conjunto de hiperparâmetros, resumidos na
Tabela~\ref{tab:hiperparametros}. Os valores relativos ao reconhecimento do número de camisola foram
definidos por experimentação, procurando equilibrar a fiabilidade da leitura com o custo
computacional, ao passo que os limiares da atribuição de equipa e do lado do campo resultam da
recalibração descrita na Secção~\ref{sec:mapeamento}, conduzida sobre os conjuntos de teste e de
validação.


\begin{table}[htbp]
    \centering
    \caption{Hiperparâmetros dos módulos intervencionados e respetiva função.}
    \label{tab:hiperparametros}
    \begin{tabular}{lll}
        \hline
        \textbf{Parâmetro} & \textbf{Valor} & \textbf{Função} \\
        \hline
        \multicolumn{3}{l}{\emph{Reconhecimento de número de camisola (pose e recorte)}} \\
        Confiança da pose            & $0.25$          & Deteção mínima dos \emph{keypoints} \\
        Alargamento lateral da ROI   & $0.55$          & Margem lateral sobre o tronco \\
        Margem superior da ROI       & $0.10$          & Recorte acima dos ombros \\
        Fração de altura do tronco   & $0.90$          & Extensão vertical da região dorsal \\
        \hline
        \multicolumn{3}{l}{\emph{Reconhecimento de número de camisola (filtro de qualidade)}} \\
        Dimensão mínima do recorte   & $25\times30$~px & Rejeita jogadores distantes \\
        Nitidez mínima               & $40.0$          & Variância do Laplaciano \\
        \hline
        \multicolumn{3}{l}{\emph{Reconhecimento de número de camisola (leitura e votação)}} \\
        Modelo de leitura            & Qwen2.5-VL-3B (4 bits) & Inferência do número \\
        \emph{Tracklets} por chamada & $4$             & Agrupamento de chamadas ao modelo \\
        Recortes por chamada         & $20$            & Limite antes de saturar a atenção \\
        Concordância na janela       & $0.75$          & Consenso dentro de cada janela \\
        Concordância no \emph{tracklet} & $0.45$       & Consenso entre janelas, com nulos \\
        \hline
        \multicolumn{3}{l}{\emph{Atribuição de equipa (agrupamento)}} \\
        Limiar de silhueta           & $0.28$          & Abaixo, recorre à cor da camisola \\
        Concentração num dos lados   & $0.85$          & Condição do teste complementar \\
        Distância entre centros      & $2.04$          & Abaixo, separação por ReID frágil \\
        Fração mínima do grupo menor & $0.20$          & Aceitação da alternativa de cor \\
        Dimensão do descritor de cor & $27$            & Histograma HSV e cor média \\
        \hline
        \multicolumn{3}{l}{\emph{Atribuição de equipa (lado do campo)}} \\
        Maioria na votação           & $0.60$          & Fração dominante por \emph{tracklet} \\
        Diferença relativa mínima    & $0.10$          & Margem do voto relativo \\
        Distância mínima entre médias & $0.3$~m        & Aciona o desempate final \\
        \hline
    \end{tabular}
\end{table}


\subsection{Protocolo de Avaliação}

A avaliação é conduzida em dois níveis complementares, seguindo a lógica exposta na
Secção~\ref{sec:metricas}. Ao nível do sistema completo, o desempenho é medido pela métrica oficial
do desafio, o GS-HOTA, que articula a precisão da localização no campo com a correção de todos os
atributos de identidade de cada interveniente. Ao nível de cada módulo intervencionado, e para isolar
o seu contributo, mede-se separadamente a exatidão da leitura do número de camisola e a exatidão da
atribuição de equipa, por comparação direta com a anotação de referência.

A recalibração dos limiares seguiu um protocolo estrito de ausência de regressões. Cada ajuste foi
testado em simultâneo nos conjuntos de teste e de validação e só foi aceite quando corrigia casos
anteriormente falhados sem degradar nenhum vídeo já correto. Deste modo, garante-se que os ganhos
reportados no capítulo seguinte resultam de melhorias consistentes em todo o conjunto avaliado, e não
de um sobreajuste a casos particulares.
